\subsection{Composition photonique du niveau \texorpdfstring{\(A_2\)}{A2} de Bendersky--Glaisher}
\label{v:subsec:bendersky-fotonica}

La densité de nombre de la proposition~\ref{v:prop:densidades-termicas}
admet une seconde lecture exacte qui conserve sa provenance spectrale.
Soit \(\mathcal Z_3\) la fonctionnelle scalaire normalisée produite par le
lecteur spectral tritique et reconnue ultérieurement comme la fonction zêta
dans la réalisation archimédienne. Soient
\[
 H_k=\sum_{j=1}^{k}\frac1j,\qquad H_0=0,
\]
et \(\mathsf B_j\) les nombres de Bernoulli avec la convention
\(\mathsf B_1=-1/2\). Pour éviter toute collision avec la notation des réseaux,
nous définissons les niveaux de Bendersky--Glaisher au moyen de
\begin{equation}
 A_k^{\mathrm{BG}}
 =\exp\!\left(
   \frac{H_k\mathsf B_{k+1}}{k+1}-\mathcal Z_3'(-k)
 \right),\qquad k\in\mathbb N_0.
 \label{v:eq:bendersky-bg}
\end{equation}
Le terme de Bernoulli fixe la normalisation du produit régularisé ;
il n’introduit pas de constante physique. La coordonnée
\(\pi_{\mathrm{HMT}}\) qui intervient dans l’équation fonctionnelle a déjà été
exprimée par l’état enrichi APP--TRIT--TPK avant cette lecture.

\begin{proposition}[Composition photonique du niveau deux]
\label{v:prop:bendersky-fotonica}
Dans la normalisation scalaire \(\mathcal Z_3=\zeta\),
\begin{equation}
 \log A_2^{\mathrm{BG}}
 =-\mathcal Z_3'(-2)
 =\frac{\mathcal Z_3(3)}{4\pi_{\mathrm{HMT}}^2}.
 \label{v:eq:bendersky-nivel-dos}
\end{equation}
Par conséquent, dans le système de coordonnées thermiques du calendrier,
\begin{align}
 n_\gamma\ell_P^3
  &=\frac{8\log A_2^{\mathrm{BG}}}{(\log3)^3}
   =-\frac{8\mathcal Z_3'(-2)}{(\log3)^3},
   \label{v:eq:numero-fotones-bendersky}\\
 \frac{u_\gamma}{n_\gamma k_B\Theta_{108}}
  &=\frac{\pi_{\mathrm{HMT}}^2}{120\log A_2^{\mathrm{BG}}},
   \label{v:eq:energia-media-bendersky}\\
 \frac{s_\gamma}{n_\gamma k_B}
  &=\frac{\pi_{\mathrm{HMT}}^2}{90\log A_2^{\mathrm{BG}}}.
   \label{v:eq:entropia-media-bendersky}
\end{align}
\end{proposition}

\begin{proof}
La forme complétée satisfait
\[
 \pi_{\mathrm{HMT}}^{-s/2}\Gamma(s/2)\mathcal Z_3(s)
 =\pi_{\mathrm{HMT}}^{-(1-s)/2}
  \Gamma((1-s)/2)\mathcal Z_3(1-s).
\]
En \(s=-2\), le zéro simple de \(\mathcal Z_3\) annule le pôle de
\(\Gamma(s/2)\). Comparer les termes constants, en utilisant
\(\Gamma(3/2)=\sqrt{\pi_{\mathrm{HMT}}}/2\) après la reconnaissance
archimédienne, donne
\[
 \mathcal Z_3'(-2)=-\frac{\mathcal Z_3(3)}{4\pi_{\mathrm{HMT}}^2}.
\]
Comme \(\mathsf B_3=0\), la définition~\eqref{v:eq:bendersky-bg}
produit la première égalité de~\eqref{v:eq:bendersky-nivel-dos}.
Substituer
\(\mathcal Z_3(3)=4\pi_{\mathrm{HMT}}^2\log A_2^{\mathrm{BG}}\)
dans~\eqref{v:eq:densidades-calendario} démontre
\eqref{v:eq:numero-fotones-bendersky}. Enfin,
\(k_B\Theta_{108}=E_P/\log3\) ; diviser les densités d’énergie et
d’entropie par la densité de nombre donne
\eqref{v:eq:energia-media-bendersky} et
\eqref{v:eq:entropia-media-bendersky}.
\end{proof}

L’évaluation \(A_2^{\mathrm{BG}}=1.0309167521973921\ldots\) sert de
test numérique des trois identités, non de donnée génératrice. Le développement
spectral de Bendersky--Glaisher contient la hiérarchie complète des dérivées
régularisées et leur correspondance avec les valeurs spectrales impaires ; la présente section
établit la composition réciproque qui utilise spécifiquement la densité
sans dimension de nombre photonique. On ne redéfinit pas \(k_B\), on n’identifie pas une densité à un
opérateur régularisé et on ne modifie aucune des densités avec
\(\zeta(3)\) déjà démontrées.
